% !TeX root = ../fundamentals.tex
% Section 2.5 -- Relevance (synthesis; NO fresh citations)
% Draft 1 (2026-07-21). Obeys WRITING_LAW.md + GLOSSARY.md. Argument-only ledger at the section end.

\section{Relevance}
\label{sec:fund:relevance}

The preceding sections are not a catalog; they are the parts of a single argument this
dissertation makes, and this section draws them together before the paper chapters put
them to work. No new literature enters here. The point is to show how the tasks, the
representations, the multi-task machinery, and the evaluation protocol combine into one
problem.

The three prediction targets are different problems, and the field has
mostly studied the hardest of them, the exact next place. Predicting the next category
and the next region is not a lesser version of that task but a different one, with its
own baselines and its own notion of a correct answer, and treating both as co-equal
end targets is the stance this work takes. That stance only pays off if the model can
represent a visit well enough to serve both. Here the representation is the lever. The
field moved from one-hot identifiers to place embeddings, and hierarchical graph
infomax gives a strong place-level vector, but any place embedding assigns a place the
same vector on every visit. A representation that cannot tell a weekday lunch from a
Saturday night out is working against both tasks at once, and the check-in level is the
response to that limit.

Multi-task learning is the mechanism that would let the two tasks share what they have
in common, but it is not a free gain. Naive hard sharing can leave a task worse than
its single-task model, and the elaborate gradient balancers proposed to prevent this
frequently do not outperform a well-tuned fixed-weight baseline. Whether joint training
helps therefore cannot be assumed; it has to be measured, against the dedicated
single-task models, under a protocol that does not flatter it. That protocol is the
last piece. User-disjoint cross-validation, macro-F1 and Acc@10 read against named
floors and the single-task ceiling, and verbs bound to paired tests and to
non-inferiority testing are what separate a real improvement from a hopeful one. The
value of the fundamentals in this chapter is that they make the central question
answerable on honest terms.

That question presses because its parts have not been brought together. The field has
a place-level representation but not a check-in-level one built for these targets; it
has multi-task learning but almost no treatment of the next region as an end target;
and it has strong evaluation practice that mobility studies do not always apply. The
dissertation takes up the three in order. It first asks whether naive multi-task
learning helps at all, given a place-level embedding and hard parameter sharing, and
reports an honest answer for that configuration. It then asks whether the
representation, rather than the architecture, is the lever, by decomposing and
enriching the input the same model receives. It finally asks what a representation
built for check-ins unlocks for a redesigned joint model, one that, by paired
superiority tests, outperforms the dedicated single-task models on the next category
everywhere it is tested and on the next region at four of six datasets, and matches the
dedicated single-task model within a two-point margin, by non-inferiority testing, at
the other two. Chapter~5 reports these tests in full. The chapters that follow answer
these three questions in turn.

% ---------------------------------------------------------------------------
% ARGUMENT LEDGER -- 2.5 (SYNTHESIS: no new citations; all claims trace to 2.1-2.4)
%  Clause 1 -> Ch.3 (CBIC): does naive MTL help, place embedding + hard sharing? (honest null)
%  Clause 2 -> Ch.4 (CoUrb): is the representation the lever, not the architecture? (decompose/enrich input)
%  Clause 3 -> Ch.5 (MobiWac, SUBMITTED/UNDER REVIEW): what does a check-in representation unlock?
%  RESULT WORDING (bound to tests, from NORTH_STAR + GLOSSARY, NOT upgraded):
%    FIX (fact F6): win verb now EXPLICITLY bound -- "by paired superiority tests, outperforms ... on the next category
%    everywhere ... and on the next region at four of six datasets"; "matches ... by non-inferiority testing" at the other two.
%    AZ/AL pair never upgraded to a win. Forward-points to Ch.5 ("Chapter 5 reports these tests in full").
%    FIX (fact F7): banned "beat" -> "do not outperform" (was "fail to beat" in the MTL synthesis sentence).
%  four-of-six + two-point margin: confirm against Ch.5 RESULTS_BOARD.md / PAPER_PLAN §3 at adaptation (fact NOTE 10).
%  NO fresh literature; NO frontier material; NO repo codenames. Section does not restate 2.1-2.4 (no "in summary").
% ---------------------------------------------------------------------------
